\begin{theorem}[Wild quadratic exact stationary assembly]\label{mod:cases-wildquadratic-errorformula}
There are separate q-free and refined stationary assemblies.

For the q-free assembly, let \(V_0\) be a
\texttt{WildQuadraticAllEvenCoefficientView}, and let \(\mathcal A\) be a
real \texttt{ExactQuadraticPhaseAssembly} for the same conductors \(t,m\).
The record \texttt{WildQuadraticActualPhaseRows} identifies the four actual
local Lamprecht packages and their literal selected stationary pairs without
mentioning a finite-field refinement.  The
\texttt{WildQuadraticAllEvenPhaseRowCompatibility} record additionally
retains the exact proof that each of those four packages is a
\texttt{LocalLamprechtPhaseData.even} row.  This is stronger than assuming
that four numerical critical phases happen to be equal to one.

The definition of an actual even Lamprecht row makes its critical factor
literally \(1\).  Consequently Lean proves from these four exact row
witnesses that
\[
 \mathcal A.\texttt{residualHasseQuotient}=1.
\]
The q-free \texttt{WildQuadraticCommonCoordinateCompatibility} identifies
the supplied \(A,u,n,z_0,z_1\) and the correction coordinate
\[
 X=\operatorname{Tr}_{K/F}(u)+n(z_0-1)+(z_1-1)
\]
with the common coefficient view.  Given the existing q-free
\texttt{ExactQuadraticPhaseAssembly.CompleteCorrectionData}, the existing
PhaseReduction theorem is applied in its exact orientation:
\[
 \operatorname{Err}_{K/F}(\chi_F)
 =\mathcal A.\texttt{residualHasseQuotient}\,
   \psi_F(-AX).
\]
Substitution of the actual-even quotient gives the structured
\texttt{WildQuadraticAllEvenErrorFormula} and the theorem
\texttt{wildQuadratic\_allEvenErrorFormula}:
\[
 \operatorname{Err}_{K/F}(\chi_F)=\psi_F(-AX).
\]
No \(\mathfrak q\), Hasse function, affine coefficient, or refinement
correction appears in this path.

For the complementary refined assembly, let \(V\) be the existing coherent
\texttt{WildQuadraticCoefficientView}, and let \(\mathcal A\) be a real
\texttt{ExactQuadraticPhaseAssembly} for the same conductors \(t,m\).
The comparison data used by Lean are explicit equalities between supplied
choices.  They exhibit the four actual local Lamprecht packages
\[
  \chi_K,\qquad \tau,\qquad \chi_F,\qquad \tau\chi_F
\]
inside the real phase data and identify their literal admissible
denominator--numerator pairs with the four selected pairs in \(V\).
For the upper and twisted packages they identify the actual critical factor
with the corresponding table phase.  For the \(\tau\)- and \(\chi_F\)-packages
they identify it with the retained quotient-source phase.  The exact actual
phase assembly carried by \(V\) then proves that the quotient of these four
factors is the coefficient-table phase ratio.  The coordinate comparison
separately identifies the supplied \(A,u,n,z_0,z_1\) data and hence the common
coordinate
\[
 X=\operatorname{Tr}_{K/F}(u)+n(z_0-1)+(z_1-1).
\]
These are compatibility hypotheses on already supplied representatives; no
stationary quotient class is eliminated to obtain a distinguished field
element.  In particular, a high or low representative can be used here only
when this literal comparison has been supplied.

The real structured result \(\mathcal A.\texttt{Result}\) first gives
\[
 \operatorname{Err}_{K/F}(\chi_F)
 =R_H\,\psi_F(-A s)\chi_F(z_0)^{-1}\tau(z_1)^{-1}.
\]
After applying the explicit coordinate equalities \(s=\operatorname{Tr}(u)\)
and the supplied identifications of \(z_0,z_1,\tau\), this becomes
\[
 \operatorname{Err}_{K/F}(\chi_F)
 =R_H\,
   \psi_F\!\left(-A\operatorname{Tr}_{K/F}(u)\right)
   \chi_F(z_0)^{-1}\tau(z_1)^{-1}.
\]
The critical quotient keeps the manuscript's orientation
\[
 R_H=\frac{H_{\chi_K}H_\tau}{H_{\chi_F}H_{\tau\chi_F}},
\]
and the coefficient comparison proves
\(R_H=\Phi_{\mathrm{table}}\).  Lean also proves that the same error term is
the quotient of the four actual complete local Lamprecht factors in this
orientation.  The result retains the exact parameter range, the nonzero
critical denominator, and the supplied real assembly's essential break-unit
identity
\[
 \mathcal A.\texttt{twistStationaryNumerator}
 =\mathcal A.\texttt{baseStationaryNumerator}
  +\mathcal A.\texttt{stationaryScale}\,
    \mathcal A.\texttt{breakUnit}.
\]
The actual conductor decompositions and parities remain in the four
existential local packages.

For each of the two possible inverse-character corrections, Lean retains an
optional proof-bearing record.  A present record contains its actual odd
conductor decomposition \(a(\theta)=2d+1\), stationary depth \(d\),
admissible denominator \(\Gamma\), stationary numerator \(\beta\), correction
\(w\), normalized class \(c\), and affine coefficient \(\lambda\), together
with the positively oriented identity
\[
 \theta(1+w)^{-1}
 =\psi_F(-\beta w/\Gamma)\,
   \mathfrak q(c)\psi_0(\lambda c).
\]
Its \(\Gamma\) and \(\beta\) are required to be exactly the selected pair of
the corresponding \(\chi_F\)- or \(\tau\)-package.  An absent record
contributes the literal factor \(1\) and carries the exact pure
linearization.  Combining only these two justified inverse-character
identities gives
\[
 \operatorname{Err}_{K/F}(\chi_F)
 =\Phi_{\mathrm{table}}\,R_\chi R_\tau\,\psi_F(-AX),
\]
where
\[
 R_\theta=\mathfrak q(c_\theta)
   \psi_0(\lambda_\theta c_\theta)
\]
when the correction is present.  Thus every inverse, sign, affine/Hasse
factor, conductor, depth, gamma, and common norm--trace correction remains
exposed before the finite table is evaluated.

Representative independence is proved for complete local Lamprecht factors.
For an odd permitted representative change the real public transport theorem
supplies a residue translation \(a\), translates the old Hasse function in
that direction, and multiplies the new elementary value by the old Hasse
value at \(a\).  Lean transports these two pieces together and proves that
\[
 \frac{L_{\chi_K}L_\tau}{L_{\chi_F}L_{\tau\chi_F}}
\]
is unchanged.  Since the real error term was identified with precisely this
quotient, the complete error expression is representative-independent.  No
isolated odd critical phase is declared invariant.

The file exports separate downstream contracts.  In all nine nonboundary
coefficient rows, the proved finite identity is
\[
 \Phi_{\mathrm{table}}R_\chi R_\tau=1,
\]
and the nonboundary contract gives
\[
 \operatorname{Err}_{K/F}(\chi_F)=\psi_F(-AX).
\]
This does not use the later theorem that the remaining elementary phase is
one.

At the odd boundary the data directly retain
\[
 \rho\ne0,\qquad r^2=\rho,\qquad 1+\rho\ne0,\qquad
 (\gamma')^2=
 \frac{\gamma+\rho\gamma_0+\rho+r}{1+\rho}.
\]
They also derive \(1+\rho^2\ne0\).  Put
\[
 A_\rho=\frac{\rho}{1+\rho},\qquad
 B=\frac{\rho^2}{1+\rho^2},\qquad
 C=\frac{\rho}{1+\rho^2},\qquad
 \gamma_1=\frac{\gamma_0+\rho\gamma+r}{1+\rho}.
\]
Given the explicit boundary-compatibility hypothesis
\(\psi_F(-AX)=\psi_0(C)\), the boundary contract retains the two selected
correction records and proves exactly
\[
 \operatorname{Err}_{K/F}(\chi_F)
 =\mathfrak q(A_\rho)\mathfrak q(B)\mathfrak q(C)\,
   \psi_0(\gamma_1A_\rho+\gamma B+\gamma_0C+C+A_\rho).
\]
No theorem from the later nonboundary, boundary, main, or final-cancellation
nodes is assumed.

\LeanFile{LanglandsFirstMainLemma/Cases/WildQuadratic/ErrorFormula.lean}
\lean{LanglandsFirstMainLemma.LocalLamprechtPhaseData.IsEven}
\lean{LanglandsFirstMainLemma.LocalLamprechtPhaseData.criticalFactor_eq_one_of_isEven}
\lean{LanglandsFirstMainLemma.WildQuadraticActualPhaseRows}
\lean{LanglandsFirstMainLemma.WildQuadraticCommonCoordinateCompatibility}
\lean{LanglandsFirstMainLemma.WildQuadraticAllEvenPhaseRowCompatibility}
\lean{LanglandsFirstMainLemma.WildQuadraticAllEvenPhaseRowCompatibility.residualHasseQuotient_eq_one}
\lean{LanglandsFirstMainLemma.WildQuadraticAllEvenErrorFormula}
\lean{LanglandsFirstMainLemma.wildQuadratic_allEvenErrorFormula}
\lean{LanglandsFirstMainLemma.wildQuadratic_errorFormula}
\lean{LanglandsFirstMainLemma.WildQuadraticPhaseRowCompatibility.toActualPhaseRows}
\lean{LanglandsFirstMainLemma.WildQuadraticCoordinateCompatibility.toCommon}
\lean{LanglandsFirstMainLemma.WildQuadraticPhaseRowCompatibility.errorTerm_eq_four_actual_completeFactors}
\lean{LanglandsFirstMainLemma.wildQuadraticCompleteErrorExpression_representative_independent}
\lean{LanglandsFirstMainLemma.LocalLamprechtPhaseData.CompleteFactorTransport.ofOdd_translation}
\lean{LanglandsFirstMainLemma.wildQuadratic_errorFormula_nonboundary}
\lean{LanglandsFirstMainLemma.wildQuadratic_errorFormula_boundary}
\leanok
\SourceText{Proposition prop:quadratic-complete-assembly.}
\uses{mod:cases-wildquadratic-coefficienttable,mod:parameters-phasereduction,mod:parameters-stationaryclassundernorm,mod:lamprecht-stationaryclasscalculus}
\FormalizationContract
\end{theorem}
